% Standalone fragment for the degree-ten X_1(2,42) argument.
% This proof uses no finite-field enumeration or computational certificate.
\begin{lemma}\label{lem:decic242-no-j0}
Let $K$ be a number field, and let $E/K$ be an elliptic curve
with $j(E)=0$. If $E(K)$ contains a subgroup isomorphic to
$C_2\times C_{42}$, then $[K:\mathbb Q]\geq12$.
\end{lemma}
\begin{proof}
Suppose that $[K:\mathbb Q]<12$ and that such a subgroup exists.
Choose $Q\in E(K)$ of order $42$. Write $E$ as $y^2=x^3+B$, where $B\in K^*$.
Full rational $2$-torsion makes the discriminant $-27B^2$ of
$x^3+B$ a square in $K$. Hence
\[
 M=\mathbb Q(\sqrt{-3})\subseteq K,
 \qquad [K:M]\leq5.
\]
Put $\zeta=(-1+\sqrt{-3})/2$. The automorphism
$(x,y)\mapsto(\zeta x,y)$ is defined over $K$, and all
endomorphisms in $\mathbb Z[\zeta]$ are therefore defined over $K$.

We first show that $E$ has good reduction at every prime of $K$
above $7$. Multiplicative reduction is impossible, since $j(E)=0$
is integral. If the reduction were additive, prime-to-$7$ torsion
would inject into the geometric component group, whose order is
at most $4$. This contradicts the point $7Q$ of order $6$.
Thus the reduction is good.

In $\mathbb Z[\zeta]$ we have
\[
 7=\lambda\overline\lambda,
 \qquad \lambda=3+\zeta,
 \qquad \overline\lambda=3+\zeta^2,
 \qquad \lambda+\overline\lambda=5.
\]
The ideals $(\lambda)$ and $(\overline\lambda)$ are distinct
primes of $M$ above $7$. Put $S=6Q$, which has order $7$.
At least one of $\lambda S$ and $\overline\lambda S$ is nonzero,
since their sum is $5S$. Choose
$\alpha\in\{\lambda,\overline\lambda\}$ such that
\[
 R=\overline\alpha S\ne0.
\]
Then $R\in E[\alpha](K)$ has order $7$.

Both primes of $M$ above $7$ have primes above them in $K$.
Hence we may choose a prime $v$ of $K$ above $(\alpha)$,
whichever nonzero component $R$ was chosen. Put $L=K_v$.
We normalize its valuation by $v_L(L^*)=\mathbb Z$, and write
$e=v_L(7)$. Since $7$ splits in $M$,
\[
 e\leq[L:\mathbb Q_7]\leq[K:M]\leq5.
\]
The isogeny $[\alpha]$ has degree $7$. Its reduction also has
degree $7$, and its differential is zero, since its multiplier
on invariant differentials is $\alpha$. Hence its reduction is
purely inseparable. It follows that $R$ reduces to the identity
and belongs to the formal group of $E/L$.
The conjugate isogeny $[\overline\alpha]$ has unit differential
at $v$. Its reduction is separable of degree $7$, so the good
reduction is ordinary.

Let $t$ be the formal parameter of $R$. Then $t\ne0$ and
$s=v_L(t)$ is a positive integer. For the ordinary formal group,
write
\[
 [7](T)=7T+a_2T^2+\cdots+a_6T^6+a_7T^7+\cdots.
\]
Here $a_i\in7\mathcal O_L$ for $2\leq i\leq6$, and
$a_7\in\mathcal O_L^*$.
Indeed, the identity $[7]^*\omega=7\omega$ for an integral
invariant differential gives $[7]^{\prime}(T)\in7\mathcal O_L[[T]]$.
This proves the divisibility of $a_2,\ldots,a_6$; the unit
condition on $a_7$ is the height-one condition.
In the equality $[7](t)=0$, the only terms which can have minimal
valuation are $7t$ and $a_7t^7$. Consequently
\[
 e+s=7s,\qquad e=6s\geq6.
\]
This contradicts $e\leq5$.
\end{proof}
